\documentclass[11pt]{article}
\usepackage{ochamath}
\subjectcolor{2F5DA8}
\setsheet{線形代数・電気系院試補充}{ケーリー・ハミルトンの定理　演習}{https://university-math-crj.pages.dev/linear-algebra/cayley-hamilton/}
\namelinetrue
\begin{document}
\makesheethead{問題}
自作問題。本文の例題とは数値や設定が異なります。条件・途中式・検算を示してください。
\problem[累乗の還元]
$A=\begin{pmatrix}3&1\\0&2\end{pmatrix}$。$A^2,A^3,A^{-1}$ を $A,I$ で表せ。
\workspace{45mm}
\problem[一般の累乗]
前問の行列で $A^k$ を $A,I$ の一次式で求めよ。
\workspace{45mm}
\newpage
\problem[重根の累乗]
$J=\begin{pmatrix}3&2\\0&3\end{pmatrix}$ の $J^k$ を求めよ。
\workspace{45mm}
\problem[3次の逆行列]
$p_A(z)=z^3-6z^2+11z-6$ のとき $A^{-1}$ を多項式で示し、検算せよ。
\workspace{45mm}
\newpage
\problem[数列]
$u_{k+2}=5u_{k+1}-6u_k,\ u_0=1,u_1=2$ を状態行列で解け。
\workspace{45mm}
\problem[一般の場合の証明]
随伴行列の係数比較によりケーリー・ハミルトンの定理の証明の骨格を述べよ。
\workspace{45mm}
\end{document}
